\section*{الفصل \ref{ch:g9:trig} --- حساب المثلثات في المثلث القائم الزاوية}

\begin{solution}{exo:g9:trig:1}
الوتر: $\sqrt{9^2 + 12^2} = \sqrt{81 + 144} = \sqrt{225} = 15$.

وضلع القائمة الآخر: $\sqrt{17^2 - 8^2} = \sqrt{289 - 64} = \sqrt{225} = 15$.
\end{solution}

\begin{solution}{exo:g9:trig:2}
$7^2 + 24^2 = 49 + 576 = 625 = 25^2$: نعم، \omterm{def:g6:shapes:triangles}{قائم الزاوية} (بعكس
فيثاغورس)، \omterm{def:g3:shapes:rightangle}{والزاوية القائمة} مقابلة للضلع $25$.

$5^2 + 6^2 = 61 \neq 64 = 8^2$: ليس \omterm{def:g6:shapes:triangles}{قائم الزاوية}.
\end{solution}

\begin{solution}{exo:g9:trig:3}
الوتر هو $[EF]$ (المقابل \omterm{def:g3:shapes:rightangle}{للزاوية القائمة} في $D$). والضلع
المقابل للزاوية $\theta = \widehat E$ هو $[DF]$؛ والضلع المجاور
هو $[DE]$. ومنه
\[
\cos\theta = \frac{DE}{EF}, \qquad
\sin\theta = \frac{DF}{EF}, \qquad
\tan\theta = \frac{DF}{DE}.
\]
\end{solution}

\begin{solution}{exo:g9:trig:4}
ضلع القائمة المقابل: $10 \sin 28^\circ \approx 10 \times 0.469 = 4.7$.
وضلع القائمة المجاور: $10 \cos 28^\circ \approx 10 \times 0.883 = 8.8$.
(والتحقّق: $4.7^2 + 8.8^2 \approx 22 + 77 \approx 100$.)
\end{solution}

\begin{solution}{exo:g9:trig:5}
$\tan\theta = \dfrac{4.5}{6} = 0.75$، إذن
$\theta = \tan^{-1}(0.75) \approx 37^\circ$.
\end{solution}

\begin{solution}{exo:g9:trig:6}
في \omterm{def:g6:shapes:triangles}{المثلث القائم الزاوية} الذي يشكّله المراقب، والنقطة التي على
الأرض تحت الطائرة، والطائرة: يكون الاِرتفاع $120$ م هو الضلع
المقابل لزاوية الاِرتفاع، والمسافة الأفقية $d$ هي المجاور.
إذن
\[
\tan 32^\circ = \frac{120}{d},
\qquad
d = \frac{120}{\tan 32^\circ} \approx \frac{120}{0.625} = 192 \text{ م}.
\]
\end{solution}

\begin{solution}{exo:g9:trig:7}
الاِرتفاع ($1.2$ م) مقابل للزاوية؛ وطول المنحدر $L$ هو
الوتر:
\[
\sin 5^\circ = \frac{1.2}{L},
\qquad
L = \frac{1.2}{\sin 5^\circ} \approx \frac{1.2}{0.0872} \approx 13.8
\text{ م}.
\]
فيجب أن يكون طول المنحدر $13.8$ م على الأقل.
\end{solution}

\begin{solution}{exo:g9:trig:8}
من $\cos^2\theta + \sin^2\theta = 1$:
$\sin^2\theta = 1 - 0.36 = 0.64$، و $\sin\theta > 0$ من أجل زاوية
حادّة، إذن $\sin\theta = 0.8$. ومنه
$\tan\theta = \dfrac{\sin\theta}{\cos\theta} = \dfrac{0.8}{0.6} =
\dfrac43$.
\end{solution}

\begin{solution}{exo:g9:trig:9}
\emph{\omterm{def:g6:shapes:triangles}{المثلث المتساوي الساقين} \omterm{def:g6:shapes:triangles}{القائم الزاوية}} ذو ضلعَي قائمة $1$:
وتره $\sqrt{1 + 1} = \sqrt2$؛ وكل زاوية حادّة فيه $45^\circ$، و
\[
\cos 45^\circ = \frac{1}{\sqrt2} = \frac{\sqrt2}{2} = \sin 45^\circ,
\qquad
\tan 45^\circ = \frac11 = 1 .
\]

\emph{\omterm{def:g6:shapes:triangles}{المثلث المتساوي الأضلاع}} ذو الضلع $1$: يقسمه اِرتفاعه إلى
مثلثين قائمَي الزاوية وترهما $1$، وقاعدتهما $\frac12$،
واِرتفاعهما $h = \sqrt{1 - \frac14} = \frac{\sqrt3}{2}$
(بفيثاغورس). والزاويتان $60^\circ$ (عند القاعدة) و $30^\circ$
(عند القمّة). وبقراءة النسب:
\[
\cos 60^\circ = \frac{1/2}{1} = \frac12, \quad
\sin 60^\circ = \frac{\sqrt3/2}{1} = \frac{\sqrt3}{2}, \quad
\cos 30^\circ = \frac{\sqrt3}{2}, \quad
\sin 30^\circ = \frac12 .
\]
\end{solution}

\begin{solution}{pb:g9:trig:1}
\textbf{1.} العلوّ فوق مستوى العين هو الضلع المقابل للزاوية
$35^\circ$، والضلع المجاور $60$ م:
$h = 60 \tan 35^\circ \approx 60 \times 0.700 = 42$ م.

\textbf{2.} $\tan\theta = 0.12$ يعطي
$\theta = \tan^{-1}(0.12) \approx 6.8^\circ$ --- وهو حادّ
بالنسبة إلى طريق، لكنّه زاوية متواضعة. والطريق بزاوية
$45^\circ$ له $\tan 45^\circ = 1$: أي اِنحدار $100\,\%$. («مئة
في المئة» بعيدة جدًّا عن العمودي --- وهو سوء فهم محبَّب عند
مصاعد التزلّج.)

\textbf{3.} في \omterm{def:g6:shapes:triangles}{مثلث قائم الزاوية} زاويتاه الحادّتان $\theta$
و $90^\circ - \theta$، يكون الضلع المقابل لإحداهما مجاورًا
للأخرى، والوتر مشترك. إذن
$\sin(90^\circ - \theta) = \frac{\text{المقابل لها}}{h} =
\frac{\text{المجاور للزاوية }\theta}{h} = \cos\theta$،
وبالتناظر $\cos(90^\circ - \theta) = \sin\theta$: \omterm{def:g9:trig:ratios}{فجيب}
\emph{التمام} هو \omterm{def:g9:trig:ratios}{الجيب} لتمام الزاوية.

\textbf{4.} الاِرتفاع: $8 \sin 60^\circ = 8 \times
\frac{\sqrt3}{2} = 4\sqrt3 \approx 6.93$ م. والقدم:
$8 \cos 60^\circ = 4$ م بالضبط.

\textbf{5.} $0.574^2 + 0.819^2 = 0.329 + 0.671 = 1.000$ (في حدود
\omterm{def:g6:decimals:rounding}{التقريب})؛ و $\left(\frac12\right)^2 +
\left(\frac{\sqrt3}{2}\right)^2 = \frac14 + \frac34 = 1$. وكل
زوج قيم يخفق في المتطابقة يفضح قراءة خاطئة أو آلة حاسبة في
النمط الخطأ --- أي كاشف أخطاء مجّاني.

\textbf{6.} $\tan 40^\circ = \dfrac hx$
و $\tan 30^\circ = \dfrac{h}{x + 200}$.

\textbf{7.} $x = \dfrac{h}{\tan 40^\circ}$
و $x + 200 = \dfrac{h}{\tan 30^\circ}$؛ وبالطرح،
$200 = h\left(\frac{1}{\tan 30^\circ} -
\frac{1}{\tan 40^\circ}\right)$، ومنه الصيغة. والأعداد:
$\frac{1}{0.577} \approx 1.732$، و $\frac{1}{0.839} \approx
1.192$، \omterm{ex:g1:subtraction:difference}{والفرق} $0.540$:
$h \approx \frac{200}{0.540} \approx 370$ م.

\textbf{8.} $x = \frac{h}{\tan 40^\circ} \approx
\frac{370}{0.839} \approx 441$ م. والتحقّق من $A$:
$\frac{h}{x + 200} \approx \frac{370}{641} \approx 0.577 =
\tan 30^\circ$: فمتّسق.

\textbf{9.} تحتاج طريقة المحطّة الواحدة إلى المسافة الأفقية إلى
النقطة \emph{الواقعة تحت القمّة مباشرةً} --- وهي غير قابلة
للقياس عبر واد أو داخل جبل. أمّا طريقة المحطّتين فتستبدل بها
مسافة تستطيع أن تمشيها بنفسك: أي $200$ م بين موضعَي نظرك أنت.

\textbf{10.} $h = \dfrac{200}{\sqrt3 - \frac{1}{\sqrt3}}
= \dfrac{200}{\frac{2}{\sqrt3}} = 100\sqrt3 \approx 173$ م ---
أي الطابق الثاني من برج إيفل، مقيسًا برؤيتين ونزهة.

\textbf{11.} خطّ النظر $d$، ونصفا \omterm{def:g4:geometry:circle}{القطر} $R$ (إلى نقطة الملامسة،
وهو عمودي على خطّ النظر) و $R + h$ (إلى العين): يعطي
فيثاغورس $d^2 + R^2 = (R + h)^2 = R^2 + 2Rh + h^2$، إذن
$d^2 = 2Rh + h^2 \approx 2Rh$ (فالعدد $h$ ضئيل أمام $R$).
وبعينين على $1.80$ م:
$d \approx \sqrt{2 \times 6.37 \times 10^6 \times 1.8}
\approx 4\,800$ م --- فأفق البحر لا يكاد يبعد خمسة
كيلومترات.

\textbf{12.} $d \approx \sqrt{2 \times 6.37 \times 10^6 \times
324} \approx 64$ كلم. وقاعدة البحّارة: $3.6\sqrt{1.8} \approx
4.8$ كلم و $3.6\sqrt{324} = 64.8$ كلم --- وكلاهما يوافق.

\textbf{13.} $\tan\theta = \frac{2}{60}$ يعطي $\theta \approx
1.9^\circ$. والزاوية $0.5^\circ$ للقمر نحو \emph{ربع} إبهام:
فالقمر يبدو هائلًا ويخفيه ظفر --- والعين منقلة رديئة.

\textbf{14.} الدرج: $\tan^{-1}\frac{17}{29} \approx 30^\circ$
--- فالدرج المريح يصعد بزاوية تكاد تكون بالضبط
$30^\circ$ من القيم الخاصّة. وبزاوية $60^\circ$ على نائمة $29$
سم، تكون القائمة $29\tan 60^\circ = 29\sqrt3 \approx
50$ سم: أي سلّم، لا درج.

\textbf{15.} $\tan 10^\circ \approx 0.18$؛
و $\tan 45^\circ = 1$؛ و $\tan 80^\circ \approx 5.7$؛
و $\tan 89^\circ \approx 57.3$. ولما $\theta \to 90^\circ$
يتضاءل الضلع المجاور حتى ينعدم بينما يصمد الضلع المقابل:
\omterm{def:g3:division:remainder}{فخارج القسمة} ينمو فوق كل حدّ. وعند $90^\circ$ بالضبط لا مثلث
ولا قيمة --- ويصعد التمثيل البياني للظلّ خطًّا مقاربًا
عموديًّا، أول كثير منها في مجلّد الثانوية.
\end{solution}
